\chapter{Геометрия спина и первое запутанное состояние}

\lectureoverview{Запутанность относится к системам, состоящим из частей, но к задаче о двух спинах нельзя по-настоящему перейти, пока не завершён аппарат одного спина. Поэтому мы начнём со спина вдоль произвольной оси, выведем вероятность перехода между двумя направлениями, истолкуем измерение как фильтр и установим, что критерием совместимости наблюдаемых служит коммутирование. Затем подсчёт показывает, что у двух спинов чистых состояний больше, чем могут дать два отдельно приготовленных спина. Эти избыточные состояния запутаны, и простейший пример уже обладает определёнными парными корреляциями, тогда как каждый электрон по отдельности полностью неполяризован.}

\section{Завершение задачи об одном спине}

С теоремой Белла можно подождать. Сама запутанность начинается с двух электронов, но вся оставшаяся структура одного электрона войдёт в каждый расчёт. Поэтому сначала доведём до конца двухсостоянийную систему, а уже потом объединим две её копии.

Выберем обычный единичный вектор в пространстве
\begin{equation}
  \hat n=(n_1,n_2,n_3),
  \qquad
  n_1^2+n_2^2+n_3^2=1.
  \label{eq:l04-unit-pointer}
\end{equation}
Метки \(1,2,3\) с тем же успехом можно назвать \(x,y,z\). Важно лишь, чтобы оси образовывали одну согласованную правую систему координат.

\begin{classroomqa}
\lecturetimestamp{00:02:33}
\audiencequestion{Нарисованный на доске координатный триэдр правый?}
\lecturerresponse{Первый рисунок левый. Если обратить вторую ось или согласованно переобозначить две оси, это исправит дело. В расчёте ниже используется только некоторая фиксированная правая конвенция; случайная ориентация рисунка не несёт никакой физики.}
\end{classroomqa}

У спина есть три наблюдаемые компоненты, представляемые матрицами Паули
\begin{equation}
  \begin{aligned}
    \sigma_1&=
    \begin{pmatrix}0&1\\1&0\end{pmatrix},
    &
    \sigma_2&=
    \begin{pmatrix}0&-i\\i&0\end{pmatrix},\\[0.5em]
    \sigma_3&=
    \begin{pmatrix}1&0\\0&-1\end{pmatrix}.
  \end{aligned}
\end{equation}
Компонента вдоль \(\hat n\) есть обычное пространственное скалярное произведение вектора-указателя с этим вектором, компоненты которого являются матрицами:
\begin{equation}
  \sigma_{\hat n}
  \equiv
  \boldsymbol{\sigma}\cdot\hat n
  =
  n_1\sigma_1+n_2\sigma_2+n_3\sigma_3.
  \label{eq:l04-sigma-dot-n}
\end{equation}
Это скалярное произведение берётся в физическом трёхмерном пространстве. Это не внутреннее произведение двух векторов состояния.

\begin{figure}[htbp]
\centering
\includegraphics[width=0.82\linewidth]{lecture_04_figure_02.png}
\caption{Пространственная компонента вдоль \(\hat n\) выражается как линейная комбинация трёх матричных наблюдаемых.}
\lectureframe{00:05:09}
\end{figure}

Подставляя матрицы Паули, получаем
\begin{equation}
  \sigma_{\hat n}
  =
  \begin{pmatrix}
    n_3 & n_1-in_2\\
    n_1+in_2 & -n_3
  \end{pmatrix}.
  \label{eq:l04-sigma-n-matrix}
\end{equation}
Удобно ввести сокращения
\begin{equation}
  n_-=n_1-in_2,
  \qquad
  n_+=n_1+in_2=n_-^*,
\end{equation}
так что
\begin{equation}
  \sigma_{\hat n}
  =
  \begin{pmatrix}
    n_3&n_-\\
    n_+&-n_3
  \end{pmatrix}.
\end{equation}
Действительные диагональные элементы и взаимно сопряжённые внедиагональные элементы делают оператор эрмитовым.

\begin{figure}[htbp]
\centering
\includegraphics[width=0.84\linewidth]{lecture_04_figure_03.png}
\caption{Наблюдаемая спиновая компонента вдоль произвольной оси в явном виде как матрица два на два.}
\lectureframe{00:07:34}
\end{figure}

С операциональной точки зрения измерить \(\sigma_{\hat n}\) значит направить сильное магнитное поле вдоль \(\hat n\) и наблюдать зависящий от спина отклик, представленный на лекции тем, испускается ли фотон. Матрица и лабораторный вопрос — это два описания одного и того же измерения.

Весь дальнейший расчёт опирается на два небольших тождества. Во-первых,
\begin{equation}
  \sigma_{\hat n}^{\,2}=I.
  \label{eq:l04-sigma-n-square}
\end{equation}
Это следует либо из умножения уравнения~\eqref{eq:l04-sigma-n-matrix} самого на себя, либо из использования
\begin{equation}
  \sigma_i^2=I,
  \qquad
  \sigma_i\sigma_j=-\sigma_j\sigma_i
  \quad(i\ne j).
\end{equation}
Следовательно, каждое собственное значение \(\lambda\) оператора \(\sigma_{\hat n}\) удовлетворяет \(\lambda^2=1\), так что возможны только результаты
\begin{equation}
  \lambda=\pm1.
\end{equation}
Во-вторых,
\begin{equation}
  n_+n_-
  =n_1^2+n_2^2
  =1-n_3^2
  =(1-n_3)(1+n_3).
  \label{eq:l04-nplus-nminus}
\end{equation}
Это не глубокая теорема, но именно это тождество держит под контролем алгебру собственных векторов.

\section{От приготовления к вероятности}

Приготовим электрон в состоянии со спином \(+1\) вдоль \(\hat n\), а затем измерим вдоль другого направления \(\hat m\). Основной вопрос для одного спина таков:
\begin{equation}
  P(m,+\mid n,+)
  =
  P(\sigma_{\hat m}=+1\mid\sigma_{\hat n}=+1).
  \label{eq:l04-central-probability}
\end{equation}
Симметрия подсказывает, что ответ должен зависеть только от угла между \(\hat n\) и \(\hat m\), и что перестановка этих двух направлений не должна его менять. Эти ожидания — полезные проверки, а не замена вероятностному постулату.

\begin{classroomqa}
\lecturetimestamp{00:14:03}
\audiencequestion{Сама квантовая механика говорит нам, что три матрицы Паули являются операторами спина, или это узнаётся из эксперимента?}
\lecturerresponse{Это отождествление является частью накопленной эмпирической структуры квантовой механики, а не результатом одного изолированного эксперимента. Повторяющиеся эксперименты со спином выявляют вероятностные закономерности, и матричный формализм был построен в работах Эйнштейна, Борна, Гейзенберга, Дирака и других, чтобы закодировать эти закономерности. Постулаты — исключительно успешные догадки, ограниченные данными.}
\end{classroomqa}

Пусть
\begin{equation}
  \sigma_{\hat n}\ket{n,+}=\ket{n,+},
  \qquad
  \sigma_{\hat m}\ket{m,+}=\ket{m,+}.
  \label{eq:l04-nm-eigenstates}
\end{equation}
Символы \(n,+\) и \(m,+\) служат метками внутри кетов. Они обозначают собственные состояния, для которых собственное значение равно \(+1\).

\begin{figure}[htbp]
\centering
\includegraphics[width=0.76\linewidth]{lecture_04_figure_04.png}
\caption{Собственные состояния с собственным значением \(+1\) для спина вдоль \(\hat n\) и \(\hat m\).}
\lectureframe{00:17:36}
\end{figure}

Первое, что нужно вычислить, — это амплитуда вероятности
\begin{equation}
  \mathcal A(m,+\mid n,+)=\braket{m,+}{n,+}.
\end{equation}
Вероятность равна квадрату её модуля:
\begin{equation}
  P(m,+\mid n,+)
  =
  \left|\braket{m,+}{n,+}\right|^2.
  \label{eq:l04-born-nm}
\end{equation}
Сначала амплитуда, потом вероятность: в этом порядке и состоит существенный квантовый шаг.

\begin{figure}[htbp]
\centering
\includegraphics[width=0.78\linewidth]{lecture_04_figure_05.png}
\caption{Перекрытие между приготовленным \(n\)-состоянием и измеряемым \(m\)-состоянием есть амплитуда вероятности.}
\lectureframe{00:19:59}
\end{figure}

\subsection{Нахождение собственного состояния}

Запишем кандидат на собственный вектор оператора \(\sigma_{\hat n}\) с собственным значением \(+1\) в виде
\begin{equation}
  \begin{pmatrix}\alpha\\\beta\end{pmatrix}.
\end{equation}
Нормировку пока можно отложить, а общая фаза не имеет физического эффекта. Поэтому вне координатной окрестности южного полюса можно выбрать \(\alpha=1\) и положить \(\beta=z\). Верхняя компонента уравнения
\begin{equation}
  \begin{pmatrix}
    n_3&n_-\\
    n_+&-n_3
  \end{pmatrix}
  \begin{pmatrix}1\\z\end{pmatrix}
  =
  \begin{pmatrix}1\\z\end{pmatrix}
\end{equation}
даёт
\begin{equation}
  n_3+n_-z=1,
  \qquad
  z=\frac{1-n_3}{n_-}.
  \label{eq:l04-z-solution}
\end{equation}
Нижняя компонента даёт то же условие после использования уравнения~\eqref{eq:l04-nplus-nminus}.

У ненормированного вектора квадрат нормы равен
\begin{align}
  \begin{pmatrix}
    1&\dfrac{1-n_3}{n_+}
  \end{pmatrix}
  \begin{pmatrix}
    1\\[3pt]\dfrac{1-n_3}{n_-}
  \end{pmatrix}
  &=
  1+\frac{(1-n_3)^2}{n_+n_-}
  \notag\\
  &=
  1+\frac{1-n_3}{1+n_3}
  =
  \frac{2}{1+n_3}.
  \label{eq:l04-unnormalized-norm}
\end{align}
Следовательно, нормированный выбор имеет вид
\begin{equation}
  \ket{n,+}
  =
  \sqrt{\frac{1+n_3}{2}}
  \begin{pmatrix}
    1\\[5pt]\dfrac{1-n_3}{n_-}
  \end{pmatrix}.
  \label{eq:l04-normalized-n-state}
\end{equation}
При \(n_3=-1\) это координатное выражение сингулярно, хотя само состояние таковым не является; собственное состояние просто равно \(\ket d=(0,1)^T\). Сингулярность относится к фазовой конвенции, а не к спину.

Заменив \(n\) на \(m\) и взяв комплексное сопряжение, получаем
\begin{equation}
  \bra{m,+}
  =
  \sqrt{\frac{1+m_3}{2}}
  \begin{pmatrix}
    1&
    \dfrac{1-m_3}{m_+}
  \end{pmatrix}.
\end{equation}
Выражение для перекрытия выглядит неприятно, а умножение его на комплексно-сопряжённое выглядит ещё хуже. Многократное использование уравнения~\eqref{eq:l04-nplus-nminus} сводит результат к
\begin{equation}
  \begin{aligned}
    P(m,+\mid n,+)
    &=\frac{1+n_1m_1+n_2m_2+n_3m_3}{2}\\
    &=\frac{1+\hat n\cdot\hat m}{2}.
  \end{aligned}
  \label{eq:l04-spin-transition-probability}
\end{equation}
Если \(\theta_{nm}\) — угол между двумя пространственными направлениями,
\begin{equation}
  P(m,+\mid n,+)
  =
  \frac{1+\cos\theta_{nm}}{2}
  =
  \cos^2\frac{\theta_{nm}}{2}.
  \label{eq:l04-half-angle-probability}
\end{equation}

\begin{classroomqa}
\lecturetimestamp{00:32:47}
\audiencequestion{Разве перекрытие не должно обращаться в нуль, поскольку два состояния ортогональны?}
\lecturerresponse{Нет. Метки \(m,+\) и \(n,+\) относятся к разным осям измерения, а не к разным собственным значениям одной фиксированной наблюдаемой. Если \(m=n\), два состояния совпадают, и их перекрытие равно единице. Ортогональными они становятся только тогда, когда физические направления противоположны, так что одно состояние есть состояние \(+1\), а другое — состояние \(-1\) одной и той же компоненты спина.}
\end{classroomqa}

\section{Что означает формула вероятности}

Уравнение~\eqref{eq:l04-spin-transition-probability} сразу проходит три проверки:
\begin{align}
  \hat m=\hat n
  &\quad\Longrightarrow\quad P=1,
  \notag\\
  \hat m=-\hat n
  &\quad\Longrightarrow\quad P=0,
  \notag\\
  \hat m\cdot\hat n=0
  &\quad\Longrightarrow\quad P=\frac12.
\end{align}
Ответ зависит только от относительного угла и симметричен относительно \(m\leftrightarrow n\). Поворот всей установки не меняет ни одного предсказания.

Таков замечательный компромисс, характерный для спина одна вторая. Вдоль любой оси наблюдаемый результат остаётся \(+1\) или \(-1\) и никогда не принимает дробного значения спина. Вместо этого плавная зависимость от поворота проявляется в вероятностях этих двух дискретных ответов.

\begin{classroomqa}
\lecturetimestamp{00:39:20}
\audiencequestion{Откуда мы знаем, что в повторных испытаниях речь идёт об одном и том же электроне, и как измеряется вероятность?}
\lecturerresponse{Можно удерживать один электрон в полости и следить за сохраняющимся зарядом по его электрическому потоку, заново приготовить этот электрон вдоль \(\hat n\) и повторять измерение вдоль \(\hat m\). Эквивалентно можно одинаково приготовить большой ансамбль электронов и измерить каждый из них один раз. В обеих процедурах, с точностью до статистической ошибки, долговременное соотношение исходов со знаком плюс и со знаком минус одинаково.}
\end{classroomqa}

\begin{classroomqa}
\lecturetimestamp{00:41:29}
\audiencequestion{Нужно ли измерять всех членов ансамбля в одном коллективном эксперименте?}
\lecturerresponse{Нет. Каждый член ансамбля приготовляется по одному и тому же протоколу и измеряется отдельно. Миллион таких испытаний образует один статистический эксперимент лишь в том смысле, что частоты их исходов оценивают одну и ту же вероятность.}
\end{classroomqa}

\begin{classroomqa}
\lecturetimestamp{00:42:23}
\audiencequestion{Если один электрон используется повторно, сколько фотонов он может испустить?}
\lecturerresponse{Никакой фиксированной пожизненной квоты нет. Каждый цикл включает внешнюю работу при включении магнитного поля. Если спин находится в ориентации с большей энергией, эта подведённая энергия может уйти в виде фотона при релаксации. Поэтому установка в принципе может повторять цикл приготовления и испускания неограниченно долго.}
\end{classroomqa}

Эту формулу можно получить очень быстро, как только установлена вращательная инвариантность: выбрать \(\hat n\) в качестве третьей оси, так что \(\ket{n,+}=(1,0)^T\), и считать вероятность по верхней компоненте \(\ket{m,+}\). Но воспользоваться этим сокращённым путём с самого начала значило бы предположить именно ту симметрию, которую явный расчёт только что продемонстрировал.

\begin{classroomqa}
\lecturetimestamp{00:43:33}
\audiencequestion{Может ли более удачный выбор базиса избавить от этих некрасивых собственных векторов для произвольной оси?}
\lecturerresponse{Да, как только в распоряжении имеются повороты и унитарные преобразования, можно совместить одну из координатных осей с одним из направлений и сделать расчёт коротким. На этом этапе явный расчёт служит определённой цели: он проверяет, что выбранные матрицы Паули действительно дают вращательно-инвариантный ответ, а не просто предполагают этот факт.}
\end{classroomqa}

\begin{classroomqa}
\lecturetimestamp{00:48:30}
\audiencequestion{В чем физическая причина того, что ответ зависит только от угла?}
\lecturerresponse{Структурный ответ относится к квантовой теории симметрии: пространственные повороты представляются унитарными преобразованиями, действующими на спиновые состояния и наблюдаемые. До этого языка мы дойдем позже. Пока зависимость от угла — это вывод, полученный вычислением с матрицами Паули.}
\end{classroomqa}

\section{Измерение — это фильтр}

Измерение не только сообщает число. Если измеряется \(\sigma_{\hat m}\) и результат равен \(+1\), электрон после этого остается в состоянии \(\ket{m,+}\). Прибор выступает и как детектор, и как устройство подготовки состояния, подобно тому как поляризатор проверяет фотон и оставляет прошедший фотон поляризованным вдоль собственной оси.

Это правило изменения состояния делает задачу одновременного измерения конкретной. Две наблюдаемые могут одновременно иметь воспроизводимые определенные значения только тогда, когда существуют состояния, являющиеся общими собственными векторами для обеих.

\begin{classroomqa}
\lecturetimestamp{00:51:29}
\audiencequestion{Может ли одно измерение одновременно определить две разные компоненты спина, такие как \(\sigma_1\) и \(\sigma_2\)?}
\lecturerresponse{Нет. Их собственные векторы различны. Измерение одной компоненты оставляет электрон в собственном состоянии этой компоненты, но это состояние не является собственным состоянием другой. Поэтому первое измерение разрушает определенность, необходимую для второго.}
\end{classroomqa}

Общее конечномерное утверждение наиболее ясно формулируется для эрмитовых операторов.

\begin{theorem}
Два эрмитовых оператора \(A\) и \(B\) одновременно диагонализуемы тогда и только тогда, когда они коммутируют:
\begin{equation}
  [A,B]\equiv AB-BA=0.
  \label{eq:l04-commutator}
\end{equation}
Эквивалентно, у них существует полный общий ортонормированный базис из собственных векторов.
\end{theorem}

Для простого направления предположим, что
\begin{equation}
  A\ket{\alpha_i}=\lambda_i\ket{\alpha_i},
  \qquad
  B\ket{\alpha_i}=\mu_i\ket{\alpha_i}
\end{equation}
для полного базиса. Тогда
\begin{align}
  AB\ket{\alpha_i}
  &=\lambda_i\mu_i\ket{\alpha_i},
  \notag\\
  BA\ket{\alpha_i}
  &=\mu_i\lambda_i\ket{\alpha_i},
\end{align}
поэтому \(AB\) и \(BA\) совпадают на каждом базисном векторе. Следовательно, \(AB=BA\). Обратно, если эрмитовы \(A\) и \(B\) коммутируют, \(B\) сохраняет каждое собственное подпространство \(A\); диагонализуя \(B\) в каждом вырожденном собственном подпространстве, получаем общий собственный базис.

\begin{figure}[htbp]
\centering
\includegraphics[width=0.76\linewidth]{lecture_04_figure_06.png}
\caption{Две наблюдаемые могут иметь одни и те же собственные векторы, но приписывать им разные собственные значения.}
\lectureframe{00:55:45}
\end{figure}

\begin{classroomqa}
\lecturetimestamp{00:53:56}
\audiencequestion{Можно ли приложить магнитные поля вдоль \(x\) и \(y\) одновременно и назвать это одновременным измерением?}
\lecturerresponse{Поля складываются в одно результирующее направление. Тогда прибор измеряет спин вдоль этой результирующей оси, возможно, под углом \(45^\circ\), а не измеряет по отдельности \(\sigma_x\) и \(\sigma_y\). Такая попытка порождает третью наблюдаемую, а не два совместимых показания.}
\end{classroomqa}

\section{Каждому чистому односпиновому состоянию соответствует ось}

Теперь обратим задачу о собственных векторах. Для произвольного нормированного спинора
\begin{equation}
  \ket\psi=
  \begin{pmatrix}\alpha\\\beta\end{pmatrix},
  \qquad
  |\alpha|^2+|\beta|^2=1,
  \label{eq:l04-general-spinor}
\end{equation}
существует ли такое пространственное направление \(\hat n\), что
\begin{equation}
  \sigma_{\hat n}\ket\psi=\ket\psi?
  \label{eq:l04-reverse-eigenproblem}
\end{equation}
Да. Сначала в лекции объясняется, почему этого следует ожидать, простым подсчетом параметров. Две комплексные компоненты содержат четыре вещественных числа. Нормировка убирает одно, а физически несущественная общая фаза — еще одно, так что остаются два. У единичной стрелки тоже есть два независимых вещественных параметра, например долгота и широта.

\begin{classroomqa}
\lecturetimestamp{01:03:52}
\audiencequestion{Как у состояния могут быть только две степени свободы, если \(\alpha\) и \(\beta\) обе являются комплексными?}
\lecturerresponse{Изначально это четыре вещественные компоненты. Условие нормировки убирает одну вещественную степень свободы, а умножение всего состояния на \(e^{i\theta}\) не меняет никаких вероятностей, так что общая фаза убирает еще одну. Остаются два физически существенных параметра, что соответствует двум углам единичной стрелки.}
\end{classroomqa}

\begin{classroomqa}
\lecturetimestamp{01:04:20}
\audiencequestion{Нужно ли нам также находить собственное значение?}
\lecturerresponse{Здесь искомое собственное значение фиксировано и равно \(+1\). Мы ищем только направление, для которого данный спинор является состоянием со спином вверх. Стрелка, направленная в противоположную сторону, описывала бы тот же луч как состояние со спином вниз с собственным значением \(-1\).}
\end{classroomqa}

Подсчет объясняет, почему такое соответствие должно существовать. Непосредственное построение это подтверждает. Определим
\begin{align}
  n_1&=\bra\psi\sigma_1\ket\psi
      =2\operatorname{Re}(\alpha^*\beta),
  \notag\\
  n_2&=\bra\psi\sigma_2\ket\psi
      =2\operatorname{Im}(\alpha^*\beta),
  \notag\\
  n_3&=\bra\psi\sigma_3\ket\psi
      =|\alpha|^2-|\beta|^2.
  \label{eq:l04-bloch-components}
\end{align}
Нормировка дает \(n_1^2+n_2^2+n_3^2=1\), а прямая подстановка дает уравнение~\eqref{eq:l04-reverse-eigenproblem}. Таким образом, каждое чистое односпиновое состояние является состоянием со спином вверх вдоль некоторой пространственной оси.

Существует и операциональная версия этого утверждения. Подготовьте много копий одного и того же неизвестного чистого состояния. Измеряйте серии вдоль пробных направлений, подстраивайте магнит и ищите направление, вдоль которого каждый результат дает спин вверх. Конечный объем данных определяет направление лишь с неопределенностью, но идеальному состоянию соответствует одна точная ось.

Вопросы — часть этого метода, а не помеха ему. Короткое отступление на лекции напоминает о том, как десятилетние дети свободно спрашивают о теории струн, дополнительных измерениях и о том, вызывают ли окружающие нейтрино головные боли. Детали вопроса могут быть неверны; сама готовность выявить точку непонимания все равно ценна.

\section{Два спина и локальные операторы}

Прежде чем объединять два спина, запишем, как матрицы Паули действуют в базисе \(\sigma_3\):
\begin{align}
  \sigma_1\ket u&=\ket d,
  &
  \sigma_1\ket d&=\ket u,
  \notag\\
  \sigma_2\ket u&=i\ket d,
  &
  \sigma_2\ket d&=-i\ket u,
  \notag\\
  \sigma_3\ket u&=\ket u,
  &
  \sigma_3\ket d&=-\ket d.
  \label{eq:l04-pauli-action-table}
\end{align}

\begin{figure}[htbp]
\centering
\includegraphics[width=0.80\linewidth]{lecture_04_figure_07.png}
\caption{Таблица действия матриц Паули представляет матрицы в виде правил действия на базисе вверх/вниз.}
\lectureframe{01:09:35}
\end{figure}

Теперь представим себе два электрона, удерживаемые в различимых местах, так что меняются только их спины. Четыре базисных состояния таковы
\begin{equation}
  \ket{uu},
  \qquad
  \ket{du},
  \qquad
  \ket{ud},
  \qquad
  \ket{dd}.
  \label{eq:l04-two-spin-basis}
\end{equation}
Первая метка относится к электрону \(1\), вторая — к электрону \(2\). Они взаимно ортонормированы, поскольку представляют различные совместные результаты двух измерений третьей компоненты спина.

\begin{figure}[htbp]
\centering
\includegraphics[width=0.82\linewidth]{lecture_04_figure_08.png}
\caption{Четыре метки базиса произведения порождают пространство состояний двух спинов.}
\lectureframe{01:11:40}
\end{figure}

Операторы, действующие на электрон \(1\), обозначим \(\sigma_i\), а действующие на электрон \(2\) — \(\tau_i\). Более явно,
\begin{equation}
  \sigma_i\equiv\sigma_i\otimes I,
  \qquad
  \tau_i\equiv I\otimes\sigma_i.
\end{equation}
Первый набор действует только на первую метку, а второй — только на вторую:
\begin{align}
  \sigma_1\ket{ud}&=\ket{dd},
  &
  \sigma_1\ket{dd}&=\ket{ud},
  \notag\\
  \sigma_3\ket{dd}&=-\ket{dd},
  &
  \tau_2\ket{ud}&=-i\ket{uu}.
  \label{eq:l04-local-operator-examples}
\end{align}
Как только действие оператора известно на четырех базисных состояниях, линейность определяет его действие на любом состоянии двух спинов.

\begin{classroomqa}
\lecturetimestamp{01:17:47}
\audiencequestion{Почему множитель \(-i\) записан вне кета? Он действует на оба электрона?}
\lecturerresponse{Символы внутри кета — это метки базисной конфигурации. Комплексный скаляр умножает весь вектор состояния, поэтому стандартная запись выносит его наружу. Здесь он возникает из действия \(\tau_2\) на вторую метку; то, что он записан снаружи, не означает, будто к обоим электронам была применена отдельная физическая операция.}
\end{classroomqa}

\begin{classroomqa}
\lecturetimestamp{01:19:44}
\audiencequestion{Означает ли умножение на \(i\) поворот электрона в обычном пространстве?}
\lecturerresponse{Само по себе — нет. Это умножение на комплексный скаляр в пространстве состояний. Пространственные повороты позже будут представляться определенными унитарными операторами; отдельный множитель \(i\) в этой таблице действия не является буквальным поворотом в обычном пространстве.}
\end{classroomqa}

\section{Шесть параметров, а не четыре}

Общее чистое состояние пары имеет вид
\begin{equation}
  \ket\Psi
  =
  c_{uu}\ket{uu}
  +c_{du}\ket{du}
  +c_{ud}\ket{ud}
  +c_{dd}\ket{dd}.
  \label{eq:l04-general-two-spin-state}
\end{equation}
Четыре комплексных коэффициента дают восемь вещественных чисел. Нормировка убирает одно, и общая фаза убирает еще одно:
\begin{equation}
  8-1-1=6.
  \label{eq:l04-two-spin-parameter-count}
\end{equation}
Для отдельно подготовленного чистого состояния электрона \(1\) нужны два вещественных параметра, и для электрона \(2\) — еще два, итого только четыре. Недостающие два параметра — не ошибка подсчета. Они указывают на существование чистых состояний пары, для которых нельзя назначить отдельные чистые состояния ее частям.

Сначала выделим состояния, которые действительно возникают при раздельной подготовке. Если
\begin{equation}
  \ket{\psi_1}=\alpha_1\ket u+\beta_1\ket d,
  \qquad
  \ket{\psi_2}=\alpha_2\ket u+\beta_2\ket d,
\end{equation}
то
\begin{align}
  \ket{\psi_1}\otimes\ket{\psi_2}
  &=
  \alpha_1\alpha_2\ket{uu}
  +\alpha_1\beta_2\ket{ud}
  \notag\\
  &\quad
  +\beta_1\alpha_2\ket{du}
  +\beta_1\beta_2\ket{dd}.
  \label{eq:l04-product-expansion}
\end{align}
Это \emph{состояние-произведение}. У каждого электрона есть своя ось поляризации, как и было бы, если бы две магнитные подготовки выполнялись независимо.

\begin{figure}[htbp]
\centering
\includegraphics[width=0.84\linewidth]{lecture_04_figure_09.png}
\caption{Произведение двух односпиновых состояний дает ограниченное четырехчленное состояние-произведение.}
\lectureframe{01:25:15}
\end{figure}

Если расположить четыре коэффициента в виде
\begin{equation}
  C=
  \begin{pmatrix}
    c_{uu}&c_{ud}\\
    c_{du}&c_{dd}
  \end{pmatrix},
\end{equation}
состояние является произведением в точности тогда, когда \(C\) имеет ранг один, или, эквивалентно,
\begin{equation}
  c_{uu}c_{dd}-c_{ud}c_{du}=0.
  \label{eq:l04-product-condition}
\end{equation}
Большинство наборов из четырех амплитуд этому условию не удовлетворяют. Они запутаны.

\section{Первое запутанное состояние}

Первый пример в лекции таков
\begin{equation}
  \ket{\psi_+}
  =
  \frac{1}{\sqrt2}
  \left(\ket{ud}+\ket{du}\right).
  \label{eq:l04-psi-plus}
\end{equation}
Его нельзя разложить в произведение состояния электрона \(1\) и состояния электрона \(2\). Если измерение вдоль третьей оси обнаруживает у электрона \(2\) спин вниз, то у электрона \(1\) спин вверх; если же оно обнаруживает у электрона \(2\) спин вверх, то у электрона \(1\) спин вниз. У пары есть определенная информация, которой ни одна из частиц по отдельности не обладает.

\begin{figure}[htbp]
\centering
\includegraphics[width=0.82\linewidth]{lecture_04_figure_10.png}
\caption{Первая запутанная суперпозиция содержит один спин вверх и один спин вниз, но заранее не приписывает ни один из этих результатов какой-либо частице.}
\lectureframe{01:29:30}
\end{figure}

Чтобы увидеть, как выглядит каждый из электронов сам по себе, вычислим три локальных средних значения. Для первой компоненты,
\begin{equation}
  \sigma_1\ket{\psi_+}
  =
  \frac{1}{\sqrt2}
  \left(\ket{dd}+\ket{uu}\right),
\end{equation}
что ортогонально \(\ket{\psi_+}\). Следовательно,
\begin{equation}
  \bra{\psi_+}\sigma_1\ket{\psi_+}=0.
\end{equation}
Аналогично,
\begin{equation}
  \sigma_2\ket{\psi_+}
  =
  \frac{i}{\sqrt2}
  \left(\ket{dd}-\ket{uu}\right),
\end{equation}
поэтому
\begin{equation}
  \bra{\psi_+}\sigma_2\ket{\psi_+}=0.
\end{equation}
Для третьей компоненты,
\begin{equation}
  \sigma_3\ket{\psi_+}
  =
  \frac{1}{\sqrt2}
  \left(\ket{ud}-\ket{du}\right),
\end{equation}
и два диагональных вклада взаимно сокращаются:
\begin{equation}
  \bra{\psi_+}\sigma_3\ket{\psi_+}
  =\frac12(1-1)=0.
\end{equation}
Итак,
\begin{equation}
  \begin{aligned}
    \bra{\psi_+}\sigma_i\ket{\psi_+}&=0,\\
    \bra{\psi_+}\tau_i\ket{\psi_+}&=0,
    \qquad i=1,2,3.
  \end{aligned}
  \label{eq:l04-zero-local-spin}
\end{equation}

\begin{figure}[htbp]
\centering
\includegraphics[width=0.84\linewidth]{lecture_04_figure_11.png}
\caption{Преобразованное состояние и сокращение \(1-1=0\) показывают, что локальные средние значения спина обращаются в нуль.}
\lectureframe{01:35:50}
\end{figure}

Для любого направления,
\begin{equation}
  \bra{\psi_+}\sigma_{\hat n}\ket{\psi_+}=0,
  \qquad
  \bra{\psi_+}\tau_{\hat n}\ket{\psi_+}=0.
  \label{eq:l04-zero-arbitrary-local-spin}
\end{equation}
Поскольку каждое локальное измерение имеет только результаты \(\pm1\), нулевое среднее означает
\begin{equation}
  P(\sigma_{\hat n}=+1)
  =
  P(\sigma_{\hat n}=-1)
  =
  \frac12,
\end{equation}
и аналогично для второго электрона. Ни одно направление само по себе не поляризует ни один из электронов. Для чистого состояния одного спина это невозможно, но для одной части запутанного чистого состояния это вполне возможно.

\begin{classroomqa}
\lecturetimestamp{01:38:23}
\audiencequestion{Являются ли локальный результат «пополам» и тот факт, что измерение одного электрона даёт нам информацию о другом, одним и тем же утверждением?}
\lecturerresponse{Нет. Для начала не учитывайте электрон \(2\): при измерении только электрона \(1\) в любом направлении получается распределение пополам. Это устанавливает отсутствие какой бы то ни было локальной поляризации. Корреляция — отдельное совместное утверждение о том, как связаны два зарегистрированных результата. Запутанность даёт обе эти особенности, но не следует использовать одну из них как замену доказательству другой.}
\end{classroomqa}

\section{Корреляция до Белла}

Корреляция сама по себе не является чем-то исключительно квантовым. Смешайте монету в 1 цент и монету в 10 центов, разнесите их, не глядя, и посмотрите на одну из них. До просмотра у каждой из них вероятность 50 на 50 оказаться монетой в 1 цент, но если здесь обнаруживается одноцентовая монета, то мы знаем, что удалённый объект — монета в 10 центов. Пара несёт определённую информацию, даже когда каждая её часть сама по себе её не несёт.

\begin{classroomqa}
\lecturetimestamp{01:41:45}
\audiencequestion{Верно ли, что вывод о монетах в 1 и 10 центов работает только тогда, когда мы знаем, что исходная пара содержала ровно по одной монете каждого типа?}
\lecturerresponse{Именно. Классический вывод опирается на надежную подготовку. Точно так же предсказание квантовой корреляции применимо только тогда, когда система из двух электронов приготовлена в указанном совместном состоянии.}
\end{classroomqa}

\begin{classroomqa}
\lecturetimestamp{01:42:09}
\audiencequestion{Для \(\ket{\psi_+}\) равна ли нулю вероятность того, что оба измерения третьей компоненты дадут вверх? Доказывает ли это, что спины противоположны вдоль любой оси?}
\lecturerresponse{Вероятность \(uu\) в базисе третьей оси равна нулю, и два результата по третьей оси противоположны. Но второе заключение слишком сильно для этого состояния со знаком плюс. Его локальный спин равен нулю во всех направлениях, тогда как совместная корреляция зависит от направления. Универсально противоположные результаты принадлежат синглету со знаком минус, введенному в конце лекции.}
\end{classroomqa}

\begin{classroomqa}
\lecturetimestamp{01:43:51}
\audiencequestion{Можно ли представить классическую пару как две монеты, спаянные так, что одна сторона одной всегда противоположна стороне другой?}
\lecturerresponse{Это возможный классический механизм фиксированной антикорреляции, если только встряхивание не нарушает эту связь. Это по-прежнему скрытая классическая конфигурация, заданная при подготовке. Аналогия намеренно предварительна; теорема Белла покажет, где всякое такое классическое объяснение терпит неудачу.}
\end{classroomqa}

\begin{classroomqa}
\lecturetimestamp{01:44:32}
\audiencequestion{Где здесь вступает в дело закон больших чисел, если один результат уже определяет результат партнера?}
\lecturerresponse{Условная достоверность — это утверждение об одном испытании, когда заданы совместное состояние и один результат. Повторяющиеся испытания всё равно нужны, чтобы экспериментально проверить, что обещанная корреляция надежна, и оценить средние значения и неидеальные вероятности. Закон больших чисел связывает эти повторяющиеся частоты с теоретическим распределением.}
\end{classroomqa}

Два наблюдателя могут быть разделены огромным расстоянием. Наблюдение одного результата меняет то, что локальный наблюдатель знает о другом результате, но не посылает удалённому наблюдателю управляемого сообщения. На этом этапе ситуация всё ещё напоминает разнесённые монеты в 1 и 10 центов.

\begin{classroomqa}
\lecturetimestamp{01:46:46}
\audiencequestion{Правильно ли думать о квантовой паре в терминах разнесенной пары монет в 1 и 10 центов?}
\lecturerresponse{Это правильное первое классическое сравнение, но не окончательное квантовое объяснение. Теорема Белла касается измеримого различия между парой, несущей заранее существующие классические свойства, и запутанной квантовой парой. Обсуждение этого различия откладывается до тех пор, пока не будут вычислены необходимые корреляции.}
\end{classroomqa}

\section{Синглет и его подготовка}

В конце выкладки на доске заменим знак плюс на минус:
\begin{equation}
  \ket{\psi_-}
  =
  \frac{1}{\sqrt2}
  \left(\ket{ud}-\ket{du}\right).
  \label{eq:l04-singlet}
\end{equation}
После этой замены локальные средние из уравнения~\eqref{eq:l04-zero-local-spin} остаются равными нулю, а результаты по третьей оси остаются противоположными. Поперечные совместные корреляции, однако, меняются. Это состояние со знаком минус — спиновый синглет, вращательно-инвариантное состояние, в котором два спина противоположны вдоль любой общей оси измерения.

\begin{figure}[htbp]
\centering
\includegraphics[width=0.84\linewidth]{lecture_04_figure_12.png}
\caption{Последняя замена на доске заменяет симметричную суперпозицию синглетом со знаком минус, сохраняя при этом нулевой локальный спин.}
\lectureframe{01:47:25}
\end{figure}

В используемой здесь упрощенной физической картине, если сблизить два локализованных магнитных момента электронов настолько, чтобы они могли взаимодействовать, то взаимная ориентация с более высокой энергией может релаксировать с испусканием фотона. При отсутствии выделенной внешней оси более низкое спиновое состояние представляется вращательно-симметричным синглетом. Точное взаимодействие и причина, по которой относительная фаза должна быть отрицательной, требуют дальнейшего расчета.

\begin{classroomqa}
\lecturetimestamp{01:49:02}
\audiencequestion{Сводится ли практическая подготовка просто к тому, чтобы сблизить два электрона и подождать?}
\lecturerresponse{В этой идеализированной модели — да: нужно дать спин-спиновому взаимодействию подействовать и допустить радиационную релаксацию в нижнее состояние. Электроны должны лишь находиться достаточно близко и достаточно долго, чтобы их магнитное взаимодействие было эффективным. Реальный протокол подготовки должен также контролировать орбитальное состояние и конкурирующие взаимодействия.}
\end{classroomqa}

\begin{classroomqa}
\lecturetimestamp{01:49:16}
\audiencequestion{Не сводится ли это просто к принципу запрета Паули?}
\lecturerresponse{Нет. Двум электронам приписаны различимые пространственные положения, так что обсуждение их спинов пока не требует антисимметризации перекрывающихся орбитальных волновых функций. Принцип запрета становится существенным, когда одинаковые электроны занимают перекрывающиеся одночастичные состояния; настоящий аргумент о запутанности уже работает для пространственно разнесённых положений.}
\end{classroomqa}

\begin{classroomqa}
\lecturetimestamp{01:50:31}
\audiencequestion{Содержал бы обычный электронный пучок такие взаимодействующие пары естественным образом?}
\lecturerresponse{Обычно нет. Пучок разрежен, и типичные электроны остаются слишком далеко друг от друга, чтобы за доступное время полёта возникло заметное спин-спиновое взаимодействие. Столкновения или специально удерживающее взаимодействие могут создавать запутанные пары, но одно лишь пребывание в одном и том же пучке этого не обеспечивает.}
\end{classroomqa}

\begin{classroomqa}
\lecturetimestamp{01:51:45}
\audiencequestion{Какое расстояние требуется для магнитного взаимодействия?}
\lecturerresponse{Универсальной резкой границы нет; решающая величина — характерное время взаимодействия. В лекции приводится грубая микроскопическая оценка порядка \(10^{-11}\,\mathrm{cm}\) для быстрого взаимодействия. При большем расстоянии энергетическое расщепление меньше, длина волны испускаемого фотона больше, и релаксация занимает гораздо больше времени. В принципе это можно компенсировать достаточно долгим ожиданием, тогда как обычный пучок предоставляет слишком мало времени.}
\end{classroomqa}

\begin{classroomqa}
\lecturetimestamp{01:53:10}
\audiencequestion{Почему синглет должен содержать знак минус, а не плюс?}
\lecturerresponse{Это не очевидно из локальной статистики «пополам», потому что эта статистика одинакова для двух относительных фаз. Ответ заключён в том, как состояние преобразуется при поворотах и как взаимодействие отбирает полный спин, равный нулю. Вывод будет в следующей лекции.}
\end{classroomqa}

Курс теперь перешел намеченный порог. У одного чистого спина всегда есть определенная ось, задающая «вверх». Чистая пара не обязана приписывать никакого чистого состояния ни одному из своих членов, и все же она может кодировать точные соотношения между их результатами. Эта дополнительная структура и есть запутанность; теорема Белла покажет, почему никакая классическая история про монеты в 1 и 10 центов не может воспроизвести ее полностью.
